\section{Two issued groups and their reinsurance share one balance sheet}
\label{sec:insurance-connected-groups}

An insurer's accounts contain several groups at once. A profitable group
cannot lend its unearned margin to an onerous group merely because both
belong to the same company. Reinsurance may recover part of the loss, but
does so through a separately measured contract. The following extension
puts those three measurements on one balance sheet and follows their cash
through the two service years. It uses the retained AASB 17 rules for an
Australian private-sector insurer's calendar 2024 and 2025 periods.

Start with cash and equity of 100. Group A is the preceding nonzero-rate
example: premium 120, claims of 44 and 48.4 at the two year ends, initial
claim present value 80, risk adjustment 10 and CSM 30. Group B receives
premium 100 and expects claims of 66 and 72.6 at the same dates. At the
same 10\% rate, B's claim present value is
\(66/1.1+72.6/1.1^2=120\). Its separately supplied risk adjustment is 10,
so it recognises a loss of 30 and has no CSM. After premium receipt its
remaining-coverage liability is 130, including a loss component of 30.
Both groups contain contracts issued at the beginning of 2024 and provide
two equal annual coverage units. Their different profitability requires
separate groups under paragraphs 16 and 47; assume no pricing-law exception
under paragraph 20. Estimating cash flows together does not authorise pooling
their CSMs \citep[paragraphs 14--28, 32--38 and 47]{aasb17}.

The current and inception discount curves remain at the assumed 10\%.
Each issued group's risk adjustment releases 5 per year, entirely in the
service result under paragraph 81. All insurance finance enters profit or
loss under paragraph 88(a). Claims are incurred and settled at each year
end, with realised claims equal to the admitted expectations on this path.
Cash is uninvested. There are no acquisition costs, investment components,
taxes, currency effects, new contracts or estimate changes. The rates,
expected claims, service units and risk adjustments are supplied conditional
measurements; these calculations do not estimate them from insurer data.

Now add treaty R, entered into before B is recognised. It covers all B
claims at 50\%, with no exclusions, limits, settlement delay or expected
non-performance. Its premium of 70 is paid immediately. Expected recoveries
are 33 and 36.3, with present value 60. A separately supplied transferred-risk
adjustment of 5 releases 2.5 each year. Its equality to half B's adjustment
is an assumption of this example, not a general quota-share valuation rule.
Without the initial loss-recovery adjustment, the net cost held in R's CSM
would be \(70-60-5=5\), and the asset after payment would be 70.

Because R was entered into in time and covers all B's contracts, paragraphs
66A and B119C--B119D recognise recovery income of \(30\times0.5=15\).
The adjustment increases R's cost CSM to \(5+15=20\), so its complete
remaining-coverage asset after payment is \(60+5+20=85\). Its loss-recovery
component of 15 is already within that asset. It is not another receivable.
Debit B's initial service loss 30 and credit its liability 30; debit R's
asset 15 and credit recovery income 15. Receive gross premiums 220 and pay
the reinsurance premium 70. Cash is 250, issued liabilities are 120 and 130,
and the reinsurance asset is 85, leaving equity 85
\citep[paragraphs 60--68, 78 and B119C--B119F]{aasb17}.

\subsection*{Financing moves the loss component as well}

B's loss component must be allocated across claim releases, risk release
and finance under paragraphs 50--52. For this unchanged group, use the
initial proportion \(\alpha=30/130=3/13\) for all three. If \(F\) is
the period's financing of B's future claim value and \(S\) is its expected
claims plus risk-adjustment release, its loss component \(K\) follows
\[
 K_{\mathrm{close}}=K_{\mathrm{open}}+\alpha F-\alpha S.
\]
This is a declared systematic allocation for these facts. It preserves
\(K=\alpha\) times B's remaining-coverage liability and reaches zero at
coverage end. New estimates, coverage or risk patterns would require their
own reassessment of the allocation; the example does not prescribe this
ratio for every insurer \citep[paragraphs 49--52 and B124]{aasb17}.

In year one, B's claim value accrues finance of 12. Its loss component
therefore increases by \(36/13\), before service releases claims 66 and
risk adjustment 5. Of that release of 71, \(213/13\) reverses the loss
and \(710/13\) is insurance revenue. The closing liability is
\(130+12-71=71\), comprising future claim value 66 and risk adjustment 5.
The loss component is \(30+36/13-213/13=213/13\), within that liability.
Debit finance expense 12 and credit B's liability; then debit the liability
71, credit revenue \(710/13\) and credit loss reversal \(213/13\).
Separately incur and pay claims 66 through the incurred-claims liability.

R's claim value accrues 6 and its CSM accrues 2. The asset gains finance
income 8. Its CSM of 22 is then allocated across two equal current and
remaining service units, releasing 11. Recovery 33, risk release 2.5 and
CSM release 11 reduce remaining coverage by 46.5; the closing asset is
\(85+8-46.5=46.5\), comprising 33, 2.5 and 11. Its loss-recovery component
tracks half B's loss component: finance adds \(18/13\), service reverses
\(213/26\), and the closing component is \(213/26\). No component is
added on top of the measured asset. In this unchanged, fully covered case,
the component reversal is represented by the scheduled cash-flow and risk
release; paragraph 66(bb) must not produce a second CSM reversal for it.

The service entry debits reinsurance premium-allocation expense
\(46.5-213/26=498/13\) and loss-recovery reversal \(213/26\), and credits
R's remaining-coverage asset 46.5. Separately recognise recovery income 33
and an incurred-recovery asset 33, then collect the cash. Thus R's service
result is \(33-46.5=-13.5\), with finance income 8 shown separately.
The component determines presentation within that service result; it does
not change its total. Paragraph 86 also prevents the reinsurance premium
allocation from reducing issued insurance revenue
\citep[paragraphs 66--67, 80--88, B72, B119 and B119F]{aasb17}.

A retains its preceding result: revenue 65.5, claims expense 44 and finance
expense 11 give profit 10.5, with remaining coverage 65.5. B has service
profit 5 and finance expense 12, giving a further loss of 7. Adding R's
loss of 5.5 makes the year-one movement in equity \(-2\), after the
initial net loss of 15. Cash is \(250-44-66+33=173\), and
\(173+46.5-65.5-71=83\) confirms closing equity.

\subsection*{Settlement and timing expose different errors}

In year two A produces profit 17.1 as before. B incurs finance expense
6.6, then releases expected claims 72.6 and risk adjustment 5. Its loss
component first increases by \(99/65\), then reverses \(1164/65\), reaching
zero. Of the total release 77.6, revenue is \(776/13\) and the remainder
is loss reversal. Paying claims 72.6 leaves B with no liability and
period loss 1.6. R accrues finance income \(3.3+1.1=4.4\), then releases
recovery 36.3, risk adjustment 2.5 and all CSM 12.1. Its asset release
is 50.9; its loss-recovery component also reaches zero. Collecting 36.3
leaves reinsurance service expense 14.6 and net expense 10.2. Total
year-two profit is \(17.1-1.6-10.2=5.3\).

The following balances are after each date's premium or claim settlements.
All incurred claims and incurred recoveries have been settled; the exact
event sequence preserves those intermediate balances before the cash entries.
\begin{center}
\begin{tabular}{lrrr}
\toprule
 & Initial receipt/payment & Year one & Year two\\
\midrule
Cash & 250 & 173 & 88.3\\
A: remaining-coverage liability & 120 & 65.5 & 0\\
B: remaining-coverage liability & 130 & 71 & 0\\
R: remaining-coverage asset & 85 & 46.5 & 0\\
Equity & 85 & 83 & 88.3\\
\bottomrule
\end{tabular}
\end{center}
The final cash of 88.3 equals initial capital 100 plus gross premiums 220,
less reinsurance premium 70 and gross claims 231, plus recoveries 69.3.
Cumulative profit is \(-15-2+5.3=-11.7\). The two issued groups and the
held treaty finish with every measurement component at zero.

Suppose instead that R is entered into immediately after B's recognition,
with no intervening service, insured event or financing and otherwise the
same prospective coverage. It is a valid different history. It cannot
recover the original recognised loss under paragraph 66A because it fails
B119C's ordering condition. R then starts with cost CSM 5 and asset 70,
with no loss-recovery component. Opening equity is 70. Its year-one closing
asset is \(33+2.5+2.75=38.25\), giving total equity 74.75. Year-two cash
and equity still finish at 88.3. The difference is the timing of recognition,
not a difference in cash. Paragraph C20A’s transition exception does not apply to this ordinary
new-business example.

Incorrectly pooling A and B produces a third history of reported amounts.
Their combined premiums 220 equal claim present value 200 plus the two
supplied risk adjustments 20. A pooled calculation would therefore show
neither CSM nor initial loss, conceal B's loss of 30, and suppress the timely
treaty's recovery adjustment of 15. It reports equity 100 initially and
91.25 after year one, although final cash is again 88.3. Those accounts
balance but violate the separate-group recognition requirements. Likewise,
adding the recovery component to the full asset overstates opening equity
by 15, while omitting financing from loss allocation makes the component
run out before coverage ends. Intermediate group identity, component
movements and recognition time are essential checks; terminal cash alone
cannot distinguish these mistakes.

Actual use also requires the treaty's boundary, collectability, measurement
inputs and eligible contract mapping. Partial coverage needs the systematic
loss allocation in B119E; a nominal quota cannot supply it. Changes in those
facts or in the estimates require new measurement events under paragraphs
66--67, rather than extending this fixed schedule unchanged. The shared
event procedure commits journal and component movements for each dated
event together, while preserving separate issued and held group identities.
\citep[paragraphs 24, 63--67, 78--86 and B119C--B119F]{aasb17}

